\section{Results on ECDFS}
\label{sec:results}

\subsection{Detection}
\label{sec:results:detection}

% --- detection numbers ---
On the $\sim$790 ECDFS tiles the detector returns $\sim$413 sources per tile at
the chosen operating point (Fig.~\ref{fig:detection}). \todo{Completeness/purity
against a reference catalog; comparison to a classical VIS baseline
(notebook 06).}

\begin{figure}
  \centering
  \includegraphics[width=\linewidth]{detection_v10_selected_threshold_overlay.png}
  \caption{Detections on a representative \textit{Euclid} VIS tile at the
    selected operating point: the classical VIS reference, the fused
    center-prediction head, and the native-resolution stem head. \todo{Tighten
    legend / crop; add a zoom-in panel near a bright star to motivate
    \S\ref{sec:results:detection}.}}
  \label{fig:detection}
\end{figure}

% --- bright star suppression ---
Detection completeness is reduced near bright stars: the source count per tile
anticorrelates with the fraction of the tile under bright extended flux
($\mathrm{corr}=-0.46$). We trace this to the per-band stem normalization
(\S\ref{sec:methods:foundation}): a saturated star produces a wide region
pinned at the clamped signal-to-noise ceiling, which dominates the group-wise
feature normalization and depresses faint-source activations. We verify this is
\emph{not} a removable smooth background (subtracting a spatially varying
background changes the count negligibly) and that it is intrinsic to the
detector rather than a noise-model artifact. \todo{Quantify against MER truth:
are the lost sources real?}

\subsection{Astrometry: per-band residual and SNR scaling}
\label{sec:results:astro}

% --- headline numbers ---
The joint-centroid astrometry reduces the pooled per-band median residual from
$\sim$50$\mas$ (raw band-vs-canonical, dominated by per-band centroid scatter)
to $\sim$11--16$\mas$ per band ($\sim$11$\mas$ pooled in the NISP and red
optical bands; Fig.~\ref{fig:astro_perband}). The residuals follow the expected
$\mathrm{FWHM}/(2.355\,\mathrm{SNR})$ centroid-noise scaling
(Fig.~\ref{fig:snr}), reaching \todo{floor at high SNR}.

\begin{figure}
  \centering
  \includegraphics[width=\linewidth]{astrometry_before_after.png}
  \caption{Median per-band astrometric offset to the \textit{Euclid} VIS
    reference, before and after correction. The raw cross-instrument offset
    (red) ranges from $\sim$50$\mas$ in the red optical bands to $130\mas$ in
    $u$; the per-object correction (blue) brings every band to
    $\sim$11--16$\mas$, with NISP reaching $11\mas$. The smooth concordance
    field alone (orange) removes little of the per-source scatter, confirming
    the residual is dominated by object-level centering, not a bulk WCS term.
    \todo{Confirm whether the per-band panel should show the joint-centroid
    instrument rather than the latent-head ablation shown here (see text).}}
  \label{fig:astro_perband}
\end{figure}

\subsection{Source-injection truth test}
\label{sec:results:injection}

% --- injection ---
To measure positional error against known truth we inject sub-pixel point
sources across all ten bands at a range of signal-to-noise levels and recover
their positions with the same pipeline. At $\mathrm{SNR}=5$ (Rubin) the
classical single-band centroid errs by $\sim$47$\mas$, the joint/canonical
position by $\sim$18--19$\mas$, approaching the VIS-only localization floor of
$\sim$17$\mas$; a position fitted jointly across all bands improves on a
VIS-only anchor by $\sim$2.4--2.5$\times$ ($16.6\rightarrow6.7\mas$ at peak
$\mathrm{SNR}=5$). \todo{Confirm numbers, add table, state caveats: idealized
Gaussian PSFs, isolated sources, injected NISP sharper than the real MER PSF.}
The test also bounds the gain: where a band is intrinsically sharper than VIS,
transferring the VIS position degrades rather than improves it
(\S\ref{sec:discussion}).

\begin{figure}
  \centering
  \includegraphics[width=\linewidth]{astrometry_truth_vs_snr.png}
  \caption{Astrometric error versus true position from controlled
    source injection into held-out tiles (Rubin bands). The classical
    single-band centroid (red) is the no-fusion baseline; the joint multi-band
    position approaches the \textit{Euclid} VIS-only localization floor (brown
    dashed), and the joint ten-band classical fit (green) is the production
    position instrument. The latent-head curves (blue/purple) overlap at the
    faint end, showing the head's $\sim$17--18$\mas$ floor is intrinsic, not
    label-limited. The dotted line marks the $\sim$4--7$\mas$ scale of the real
    concordance field (\S\ref{sec:results:field}).}
  \label{fig:snr}
\end{figure}

\subsection{Absolute validation against \textit{Gaia}}
\label{sec:results:gaia}

% --- gaia ---
Cross-matching the joint canonical positions to \textit{Gaia} gives a bulk
offset of $(+3.7, +1.9)\mas$ and a per-star scatter of $\sim$5.4$\mas$ after
bulk removal, improving to $4.4$--$4.8\mas$ for $G<18.5$. \todo{Gaia DR,
matching radius, N, magnitude-binned table, spatial structure of the offset.}

\subsection{A coherent cross-survey concordance field}
\label{sec:results:field}

% --- the field ---
After removing the bulk offset, the per-band residuals between the Rubin and
\textit{Euclid} astrometric solutions are not spatially random: they trace a
smooth, coherent field across the survey footprint. Decomposing the matched
source offsets with a multi-scale field solver, we isolate a degree-scale
component of amplitude $\sim$4--7$\mas$ that varies slowly across the field
(Fig.~\ref{fig:field}). This is the residual mismatch between the two
instruments' independently solved World Coordinate Systems --- precisely the
unmodeled structure that classical per-instrument pipelines, which calibrate
each survey's astrometry separately, leave in place
(\S\ref{sec:intro}, SITCOMTN-159).

% --- why we trust it ---
We establish that this degree-scale field is a real measurement and not solver
noise: it is recovered consistently across independent fits, across the choice
of duplicate handling, across reference catalogs, and --- importantly --- across
\jaisp{} foundation-model versions, with the smoothed field agreeing at Pearson
$r$ up to $0.92$. \todo{State the smoothing scale ($\sim$150\arcsec) and the
cross-comparison numbers; reference the controlled refit audit.}

% --- honest bound ---
We are deliberate about what is and is not established. Only the degree-scale
component is reproducible; sub-degree structure in any single solver product
sits below the solver's run-to-run repeatability floor ($\sim$5$\mas$) and is
not interpretable. We therefore report the field as a characterized
\emph{measurement} of the cross-survey WCS residual, not as a production
correction applied to the catalog.

\begin{figure}
  \centering
  \includegraphics[width=\linewidth]{hgp_lattice_smoothed_vs_pinn.png}
  \caption{The degree-scale cross-survey concordance field (\textit{i} band),
    recovered two independent ways from the matched per-band residuals: a
    lattice-smoothed hierarchical-GP fit (left) and a physics-informed
    smooth-field fit (right). Both recover the same $\sim$4--7$\mas$ large-scale
    topology; their agreement at this scale --- and across foundation-model
    versions --- is the evidence that the field is a property of the data rather
    than solver noise. \todo{Optionally add the cross-version overlay /
    correlation panel.}}
  \label{fig:field}
\end{figure}
